\documentclass[11pt]{article}

\usepackage[a4paper,margin=1in]{geometry}
\usepackage{amsmath,amssymb}
\usepackage{hyperref}

\hypersetup{
    colorlinks=true,
    linkcolor=blue,
    citecolor=blue,
    urlcolor=blue
}

\title{Why Displacement Interpolation Gives Wasserstein Geodesics}
\author{}
\date{}

\begin{document}

\maketitle

Assume $\mu,\nu\in\mathcal P_2(\mathbb R^d)$ and that the optimal transport from
$\mu$ to $\nu$ is induced by a map
\[
T:\mathbb R^d\to\mathbb R^d,
\qquad
T_\#\mu=\nu.
\]
For the quadratic cost, define
\[
T_t(x)=(1-t)x+tT(x),
\qquad t\in[0,1],
\]
and
\[
\mu_t=(T_t)_\#\mu.
\]
This curve moves each particle along the straight line from $x$ to $T(x)$.

\section*{Why This Is a Geodesic}

For $0\le s<t\le 1$, the map sending $T_s(x)$ to $T_t(x)$ moves the same
particle from its position at time $s$ to its position at time $t$. The distance
traveled by that particle is
\[
|T_t(x)-T_s(x)|
=
|(t-s)(T(x)-x)|
=
|t-s|\,|T(x)-x|.
\]
Therefore the quadratic transport cost between $\mu_s$ and $\mu_t$ is bounded by
\[
W_2^2(\mu_s,\mu_t)
\le
|t-s|^2
\int_{\mathbb R^d}|T(x)-x|^2\,d\mu(x).
\]
Since $T$ is optimal from $\mu$ to $\nu$,
\[
\int_{\mathbb R^d}|T(x)-x|^2\,d\mu(x)
=
W_2^2(\mu,\nu).
\]
Thus
\[
W_2(\mu_s,\mu_t)
\le
|t-s|\,W_2(\mu,\nu).
\]

The reverse inequality follows from the triangle inequality:
\[
W_2(\mu,\nu)
\le
W_2(\mu,\mu_s)
+
W_2(\mu_s,\mu_t)
+
W_2(\mu_t,\nu).
\]
Using the same upper-bound argument on the intervals $[0,s]$ and $[t,1]$ gives
\[
W_2(\mu,\mu_s)\le sW_2(\mu,\nu),
\qquad
W_2(\mu_t,\nu)\le (1-t)W_2(\mu,\nu).
\]
Combining these inequalities forces
\[
W_2(\mu_s,\mu_t)
\ge
|t-s|\,W_2(\mu,\nu).
\]
Together,
\[
W_2(\mu_s,\mu_t)
=
|t-s|\,W_2(\mu,\nu).
\]
This is exactly the constant-speed geodesic condition.

\section*{Interpretation}

Euclidean geodesics move points along straight lines. Wasserstein geodesics move
mass along optimal straight-line particle trajectories. This is why the curve
$\mu_t=(T_t)_\#\mu$ is called displacement interpolation.

\end{document}
